\documentclass[11pt,reqno]{amsart}
\usepackage[localtheorems]{raeez-math-template}

% ================================================================
% Packages
% ================================================================
\usepackage{array}

% ================================================================
% Theorem environments
% ================================================================
\newtheorem{theorem}{Theorem}[section]
\newtheorem{proposition}[theorem]{Proposition}
\newtheorem{lemma}[theorem]{Lemma}
\newtheorem{corollary}[theorem]{Corollary}
\theoremstyle{definition}
\newtheorem{definition}[theorem]{Definition}
\newtheorem{construction}[theorem]{Construction}
\newtheorem{example}[theorem]{Example}
\theoremstyle{remark}
\newtheorem{remark}[theorem]{Remark}

% ================================================================
% Macros
% ================================================================
\newcommand{\cA}{\mathcal{A}}
\newcommand{\cH}{\mathcal{H}}
\newcommand{\cM}{\mathcal{M}}
\newcommand{\cO}{\mathcal{O}}
\newcommand{\cL}{\mathcal{L}}
\newcommand{\cW}{\mathcal{W}}
\newcommand{\Ainf}{A_{\infty}}
\newcommand{\SCchtop}{\mathsf{SC}^{\mathrm{ch,top}}}
\newcommand{\HH}{\mathrm{HH}}
\newcommand{\RR}{\mathbb{R}}
\newcommand{\Ran}{\mathrm{Ran}}
\newcommand{\MC}{\mathrm{MC}}
\newcommand{\Sym}{\mathrm{Sym}}
\newcommand{\Ass}{\mathrm{Ass}}
\newcommand{\Com}{\mathrm{Com}}
\newcommand{\Lie}{\mathrm{Lie}}
\newcommand{\Hom}{\mathrm{Hom}}
\newcommand{\Res}{\mathrm{Res}}
\newcommand{\Aut}{\mathrm{Aut}}
\newcommand{\End}{\mathrm{End}}
\newcommand{\fg}{\mathfrak{g}}
\newcommand{\gmod}{\mathfrak{g}^{\mathrm{mod}}_{\cA}}
\newcommand{\gEone}{\mathfrak{g}^{E_{1}}_{\cA}}
\newcommand{\Barch}{\overline{B}^{\mathrm{ch}}}
\newcommand{\BarchFG}{\overline{B}^{\mathrm{FG}}}
\newcommand{\Barord}{\overline{B}^{\mathrm{ord}}}
\newcommand{\BarSig}{\overline{B}^{\Sigma}}
\newcommand{\Cobar}{\Omega^{\mathrm{ch}}}
\newcommand{\Vir}{\mathrm{Vir}}
\newcommand{\KM}{\mathrm{KM}}
\newcommand{\CC}{\mathbb{C}}
\newcommand{\op}{\mathrm{op}}
\newcommand{\FM}{\overline{\mathrm{FM}}}
\newcommand{\av}{\operatorname{av}}
\newcommand{\GRT}{\mathrm{GRT}}
\newcommand{\id}{\mathrm{id}}

\DeclareMathOperator{\rk}{rk}
\DeclareMathOperator{\depth}{depth}
\DeclareMathOperator{\ch}{ch}

\numberwithin{equation}{section}

% ================================================================
\begin{document}

\title[$E_{1}$ Primacy in Modular Koszul Duality]
{$E_{1}$ Primacy in Modular Koszul Duality:\\
The Ordered Bar Complex, the Averaging Map,\\
and Five Invariants Surviving a $\Sigma_{n}$-Quotient}

\author{Raeez Lorgat}
\address{Department of Mathematics, Harvard University,
 Cambridge, MA 02138, USA}
\email{rlorgat@math.harvard.edu}

\date{}

\begin{abstract}
For a cyclic chiral algebra $\cA$, the ordered bar complex
$\Barord(\cA) = T^{c}(s^{-1}\bar\cA)$ with deconcatenation coproduct
is the level-$2$ chain object before averaging. The symmetric bar
$\BarSig(\cA)$, which underlies the classical Francis--Gaitsgory
factorization picture, is its $R$-twisted coinvariant image when
descent exists; in the pole-free locus this is the ordinary
$\Sigma_n$-coinvariant image under an admissible averaging
comparison $\av \colon \gEone \to \gmod$. The scalar modular
characteristic $\kappa(\cA)$
is recovered at degree $2$ from a meromorphic $r$-matrix $r(z)$ that
lives natively on the ordered side: for abelian families the average
is $\kappa$, while for non-abelian affine Kac--Moody it is the
double-pole term $\kappa_{\mathrm{dp}}$, and the full $\kappa$
adds the Sugawara shift $\dim(\fg)/2$. The Drinfeld associator has
a degree-$3$ coinvariant shadow, while the full associator remains
ordered data. The five foundational theorems A--D and H of modular
Koszul duality have ordered formulations on the Koszul/descent
locus; their $\Sigma_n$-shadows recover the symmetric theorems
where the descent comparison is available. Information lost under
$\av$ includes the full braiding, associator-dependent choices of
splitting, and cross-degree terms invisible to scalar invariants.
Three distinct coalgebra structures feature and must not be
conflated: the Harrison coLie structure with $n!/n$ terms at degree
$n$, the coshuffle structure with $2^{n}$ terms, and the
deconcatenation structure with $n+1$ terms. Only deconcatenation
equips the ordered bar. There are three tiers of $r$-matrix input:
the pole-free tier with Koszul-signed $R=\tau$, the vertex-algebra
tier where $R$ is read from the local OPE, and the genuinely
$E_1$ tier containing Yangians and Etingof--Kazhdan quantum vertex
algebras.

The dimensional uplift from a $2$-real-dimensional chiral algebra
on a curve to a $3$-real-dimensional holomorphic--topological
theory is governed by the chiral Deligne--Tamarkin theorem
(Theorem~\ref{thm:swiss-cheese-chiral-DT-N3} below): the universal
local Swiss-cheese pair
$(Z^{\mathrm{der}}_{\mathrm{ch}}(\cA),\, \cA)$ is terminal in the
category of two-coloured $\SCchtop$-pairs over $\cA$, the open
colour $\cA$ is acted on by the closed colour through the chiral
Swiss-cheese mixed operations, and the closed sector is not acted
on by the open sector. The bar coalgebra $\Barord(\cA)$ is a
chain-level computational model that resolves $\cA$ for the purpose
of computing $Z^{\mathrm{der}}_{\mathrm{ch}}(\cA)$ via
$\mathrm{Hom}_{\mathrm{ch}}(\Barord(\cA),\, \cA)$; it does not by
itself supply the closed colour or the dimensional uplift.
\end{abstract}

\subjclass[2020]{17B69, 18M70, 18M85, 81T40}

\keywords{Ordered bar complex, $E_{1}$-algebras, chiral algebras,
Koszul duality, $R$-matrix, Drinfeld associator,
modular operad, averaging, Grothendieck--Teichm\"uller group}

\maketitle

% ================================================================
\section{Introduction}
\label{sec:introduction}

The bar complex attached to an augmented associative algebra
admits two very different presentations. In the first, one takes
the reduced tensor coalgebra $T^{c}(s^{-1}\bar A)$ with its
deconcatenation coproduct; the bar differential is built from the
associative product in consecutive slots. In the second, one passes
to $\Sigma_{n}$-coinvariants and works with the symmetric bar
complex, which is the natural home for commutative or
$E_{\infty}$-algebras and for the Francis--Gaitsgory chiral bar.
The ordered bar complex is the level-$2$ chain object before
averaging; the symmetric bar is obtained from it by descent and coinvariants. The
averaging map is $R$-twisted when the spectral kernel is nontrivial
and ordinary only in the pole-free locus. Under an admissible
descent datum it gives a dg Lie comparison from the ordered
convolution algebra to the symmetric one. The five foundational
theorems of modular Koszul duality (Theorems A through D, together
with the chiral Hochschild theorem H) have ordered versions on the
Koszul/descent locus, and each symmetric theorem is the
$\Sigma_{n}$-coinvariant shadow of its ordered counterpart. The
information lost under $\av$ is geometric, arithmetic, and
physical: the spectral $r$-matrix, the KZ connection, the
Drinfeld associator, and the full braiding of tensor categories
of modules all live on the ordered side.

Three principles motivate the inversion. The first is operadic:
Stasheff's associahedra $K_{n}$ (cf.~\cite{Stasheff63}) carry
more information than the symmetric groups $\Sigma_{n}$, and the
Getzler--Jones~\cite{GetzlerJones94} theory of $E_{n}$-operads
makes $E_{1}$ the lowest rung of a ladder whose every step is
non-commutative. The second is representation-theoretic: the
classical $r$-matrix of an affine Kac--Moody algebra at level $k$
is the content of the degree-$2$ genus-$0$ Maurer--Cartan element
on the ordered side, and its $\Sigma_{2}$-average recovers the
degree-$2$ scalar shadow: for Heisenberg this is $\kappa$, while
for non-abelian affine Kac--Moody it is
$\kappa_{\mathrm{dp}} = k\dim(\fg)/(2h^{\vee})$, with the full
$\kappa$ obtained after adding the Sugawara shift $\dim(\fg)/2$.
The third is operadic--holographic: the chiral Deligne--Tamarkin
theorem (Theorem~\ref{thm:swiss-cheese-chiral-DT-N3} below)
constructs the universal Swiss-cheese pair
$(Z^{\mathrm{der}}_{\mathrm{ch}}(\cA),\, \cA)$ as terminal in the
category of two-coloured $\SCchtop$-pairs over $\cA$. The closed
colour $Z^{\mathrm{der}}_{\mathrm{ch}}(\cA)$ acts on the open
colour $\cA$ through the chiral Swiss-cheese mixed operations; the
open colour does not act back on the closed colour. The bar
coalgebra $\Barord(\cA)$ is a chain-level computational model used
to compute and present this Swiss-cheese pair, not the source of
the open--closed transition. The averaging map sends the ordered
structure to its $\Sigma_n$-coinvariant image \emph{within} the
open colour; the closed--open Swiss-cheese relation requires the
two-coloured operad and is invisible to averaging on a single
colour.

The conventions are those of the monograph \cite{Lorgat26I},
referred to below as Volume I. OPE-mode notation is used throughout;
conformal
weights are denoted $h$ and desuspended degrees $|s^{-1}v| = |v|-1$
(the bar complex lowers degree). The grading is cohomological.

\subsection{Notation and conventions}

Let $k$ be a field of characteristic zero, $X$ a smooth algebraic
curve over $k$, and $\cA$ an augmented cyclic chiral algebra on
$X$ in the sense of Beilinson--Drinfeld. Write $\bar\cA$ for
the augmentation ideal $\ker(\varepsilon \colon \cA \to k)$. All
bar constructions are defined on $\bar\cA$, not on $\cA$ itself.
Symmetric groups $\Sigma_{n}$ act on tensor powers in the
standard way. Over characteristic zero, the finite group algebra
$k[\Sigma_n]$ is semisimple, so invariants and coinvariants are
exact and identified by the Reynolds idempotent. Coinvariants are
denoted with a lower subscript, $(-)_{\Sigma_{n}}$. The ordered and symmetric
Fulton--MacPherson compactifications are written
$\FM_{n}^{\mathrm{ord}}(\CC)$ and $\FM_{n}(\CC)$ respectively;
the former is the Ran-space blow-up along ordered collision
diagonals.

% ================================================================
\section{Three coalgebra structures}
\label{sec:three-coalgebras}

A sum of tensor powers carries three distinct cooperadic
structures. They differ as natural transformations and by their
degree-$n$ coproduct term counts. Only deconcatenation equips the
ordered bar complex.

\begin{definition}[The three coalgebra structures on $T(V)$]
\label{def:three-coalgebras}
Let $V$ be a graded vector space and $T(V) = \bigoplus_{n\geq 0}
V^{\otimes n}$ its tensor space. Three cooperad structures on
$T(V)$ (or on the ordered/symmetric variants) arise naturally:
\begin{enumerate}[label=\textup{(\roman*)}]
\item The \emph{Harrison coLie} (or co-Lie) cooperad
 $\Lie^{c}$, whose degree-$n$ component has dimension $(n-1)!$ and
 whose coproduct decomposes a Lie word into pairs of sub-Lie-words
 using the Harrison differential.
\item The \emph{coshuffle cooperad} $\Sym^{c}$, whose degree-$n$
 coproduct produces $2^{n}$ summands by the shuffle formula
 \[
 \Delta_{\mathrm{sh}}(v_{1}\cdots v_{n})
 = \sum_{I\sqcup J = \{1,\ldots,n\}}
 \varepsilon(I,J)\,
 (v_{I}) \otimes (v_{J}),
 \]
 with the sum running over all $2^{n}$ bipartitions and
 $\varepsilon(I,J)$ the Koszul sign.
\item The \emph{deconcatenation cooperad} $T^{c}$, whose degree-$n$
 coproduct produces $n+1$ summands by splitting the word at a
 single position:
 \[
 \Delta(v_{1}\cdots v_{n})
 = \sum_{i=0}^{n}
 (v_{1}\cdots v_{i}) \otimes (v_{i+1}\cdots v_{n}).
 \]
\end{enumerate}
\end{definition}

The three structures are not interconvertible without loss of
information. At degree $3$ the Harrison differential sees
$2$ terms, the coshuffle sees $8$, and deconcatenation sees $4$.
Only the last is ordered and strictly associative at the
coalgebra level.

\begin{proposition}[Coshuffle is not deconcatenation]
\label{prop:coshuffle-not-deconcat}
Let $V$ be nonzero. The coshuffle coproduct on $\Sym(V)$ and the
deconcatenation coproduct on $T^{c}(V)$ are distinct cooperadic
structures: they have different degree-$n$ term counts
($2^{n}$ vs.\ $n+1$), different symmetry properties
(the coshuffle is cocommutative, deconcatenation is not), and
different primitive subspaces.
\end{proposition}

\begin{proof}
Cocommutativity of the coshuffle follows from the involution
$I \leftrightarrow J$ on bipartitions. Deconcatenation does not
admit such an involution: swapping $(v_{1}\cdots v_{i})$ with
$(v_{i+1}\cdots v_{n})$ does not preserve the ordered word. The
term counts are elementary.
\end{proof}

The three cooperads sit in a commutative square of natural
transformations whose top arrow is the symmetrization and whose
bottom arrow is the Harrison projection, but they are not
canonically isomorphic. Conflating coshuffle with deconcatenation
is a standard error in the chiral-algebra literature and produces
spurious identifications between factorization coproducts and
bar coproducts.

\begin{remark}[Where each structure appears]
\label{rem:where-three-appear}
The Harrison coLie cooperad $\Lie^{c}$ underlies the chiral Lie
bar complex used by Francis--Gaitsgory, where only the zeroth
product of the OPE is visible. The coshuffle cooperad $\Sym^{c}$
underlies the symmetric bar $\BarSig(\cA)$ on the Ran space,
where factorization is implemented by the topological
coproduct on unordered configuration space. The deconcatenation
cooperad $T^{c}$ underlies the ordered bar $\Barord(\cA)$, where
the coproduct splits a word along a single internal edge and
remembers the ordering on the collision points. The averaging
map $\av$ of Section~\ref{sec:av-map} descends from $T^{c}$ to
$\Sym^{c}$ by external $\Sigma_{n}$-coinvariants; it does not
factor through Harrison unless one further restricts to the
Eulerian weight-$1$ summand.
\end{remark}

% ================================================================
\section{The ordered bar complex \texorpdfstring{$\Barord(\cA)$}{Bord(A)}}
\label{sec:ordered-bar}

For a cyclic chiral algebra, the ordered bar complex is the
deconcatenation coalgebra on the ordered configuration space.

\begin{definition}[The ordered bar complex]
\label{def:ordered-bar}
Let $\cA$ be a strongly admissible augmented $E_{1}$-chiral
algebra on $X$, with augmentation ideal $\bar\cA = \ker\varepsilon$.
The \emph{ordered bar complex} is the graded vector space
\begin{equation}
\label{eq:ordered-bar-definition}
\Barord(\cA)
\;:=\;
T^{c}\bigl(s^{-1}\bar\cA\bigr)
\;=\;
\bigoplus_{n \geq 0} \bigl(s^{-1}\bar\cA\bigr)^{\otimes n},
\end{equation}
equipped with the deconcatenation coproduct
$\Delta([a_{1}|\cdots|a_{n}])
= \sum_{i=0}^{n}
 [a_{1}|\cdots|a_{i}] \otimes [a_{i+1}|\cdots|a_{n}]$
and the bar differential
$d_{\mathrm{bar}}$ obtained by extracting OPE collision residues
at consecutive pairs of points in the ordered configuration
space. The augmentation ideal $\bar\cA$ is essential here: using
$\cA$ in place of $\bar\cA$ in \eqref{eq:ordered-bar-definition}
would include the unit and break the construction.
\end{definition}

The desuspension $s^{-1}$ lowers the degree by one, so
$|s^{-1}a| = |a|-1$. The bar differential lowers tensor degree by
one; the coproduct is a chain map with respect to
$d_{\mathrm{bar}}$ by associativity of the chiral product in
consecutive slots. Both the differential and the coproduct are
built from the \emph{ordered} collision pattern on
$\FM_{n}^{\mathrm{ord}}(\CC)$; the form $\eta = d\log(z_{i}-z_{j})$
is the universal Arnold propagator (weight one, not weight $h$).

\begin{theorem}[Ordered bar differential squared]
\label{thm:ordered-d-squared}
The ordered bar differential satisfies $d_{\mathrm{bar}}^{2} = 0$.
\end{theorem}

\begin{proof}
The argument is a codimension-two face pairing on
$\FM_{n}^{\mathrm{ord}}(\CC)$ with the Arnold relation supplying
the required cancellation between the two ways of contracting an
adjacent pair of collisions. The full face-pairing argument is
the ordered Fulton--MacPherson compactification calculation of
\cite[Chapter~3]{Lorgat26I}.
\end{proof}

\begin{remark}[Three bar complexes, one map]
\label{rem:three-bars}
Let $\BarchFG(\cA)$ denote the Francis--Gaitsgory chiral Lie bar,
which uses only the zeroth OPE product $a_{(0)}b$; let
$\BarSig(\cA) = \Barch(\cA)$ denote the symmetric bar, which uses
all OPE products and takes $\Sigma_{n}$-coinvariants (the bar
complex of Volume I, Theorem A); and let $\Barord(\cA)$ denote
the ordered bar of \eqref{eq:ordered-bar-definition}. There is a
natural surjection
\begin{equation}
\label{eq:three-bars}
\Barord(\cA) \twoheadrightarrow \BarSig(\cA),
\qquad
\mathrm{gr}\,\BarSig(\cA) \twoheadrightarrow \BarchFG(\cA),
\end{equation}
where the first arrow is $\Sigma_{n}$-coinvariant projection
and the second is associated graded along the OPE filtration.
These three complexes must not be conflated: the
Francis--Gaitsgory complex sees only the Lie part, the symmetric
bar sees the full symmetric OPE, and the ordered bar sees the
braiding as well.
\end{remark}

\begin{example}[Heisenberg ordered bar]
\label{ex:heisenberg-ordered-bar}
Let $\cH_{k}$ be the Heisenberg chiral algebra at level $k$,
generated by a single field $J$ of conformal weight $h=1$. The
desuspended generator has degree $|s^{-1}J| = 0$. At degree $2$
the ordered bar component is
\[
(s^{-1}\bar\cH_{k})^{\otimes 2}
\cong k \cdot [J|J],
\]
a one-dimensional space. The bar differential extracts the
collision residue of $J(z_{1})J(z_{2}) \sim k/(z_{1}-z_{2})^{2}$
against the Arnold form $\eta = d\log(z_{1}-z_{2})$, producing
a single pole $k/(z_{1}-z_{2})$.
\end{example}

\begin{remark}[Heisenberg flat and curved $R$-data]
\label{rem:heisenberg-flat-curved-rdatum}
The Heisenberg algebra $\cH_{k}$ has a flat presentation
$r(z)=k/z$ and a curved presentation $R(z)=\exp(k\hbar/z)$.
They are related by $R(z)=\exp(\hbar r(z))$, so $r(z)$ is the
classical coefficient of $\log R$. Since $r(z)$ is central in the
two tensor factors, the exponentiated form satisfies the quantum
Yang--Baxter equation order by order in $\hbar$.
\end{remark}

% ================================================================
\section{The averaging map \texorpdfstring{$\av$}{av}}
\label{sec:av-map}

The averaging comparison realises the symmetric bar as the
$\Sigma_{n}$-coinvariant shadow of the ordered bar after a descent
datum has been chosen. Its domain is the $E_{1}$ convolution dg Lie
algebra, its codomain is the modular convolution dg Lie algebra of
Volume I, and its kernel records the information lost in the passage
from the open to the closed sector.

\begin{definition}[The $E_{1}$ modular convolution dg Lie algebra]
\label{def:e1-convolution}
Let $\cA$ be a cyclic $E_{1}$-chiral algebra and let $\cM_{\Ass}$
denote the ribbon modular operad with $\cM_{\Ass}(g,n)
= C_{*}(\overline{\cM}_{g,n}^{\,\mathrm{rib}})$, the chain complex
of the moduli of stable Riemann surfaces with $n$ parametrised
boundary circles. Define
\begin{equation}
\label{eq:e1-convolution-hom}
\gEone
\;:=\;
\prod_{\substack{g,n \\ 2g-2+n > 0}}
\Hom\!\bigl(\cM_{\Ass}(g,n),\,\End_{\cA}(n)\bigr).
\end{equation}
The Hom is taken \emph{without} $\Sigma_{n}$-equivariance: this
is the structural feature distinguishing $\gEone$ from the
modular convolution algebra $\gmod$ of Volume I, in which the
Hom is $\Sigma_{n}$-equivariant and the source cooperad is of
$\Sym^{c}$-type; the ordered source is of $T^{c}$-type. The dg Lie structure
on $\gEone$ is inherited from the Feynman transform
$F\Ass$: the differential contracts internal edges of ribbon
graphs, the bracket composes ribbon graphs.
\end{definition}

\begin{definition}[Admissible averaging datum and averaging map]
\label{def:averaging-map}
An \emph{admissible averaging datum} is a strong-unitary
$R$-twisted $\Sigma_n$-descent action $\rho_R$ on the ordered bar
local systems, satisfying the Yang--Baxter relation,
$R_{12}(z)R_{21}(-z)=\id$, and Fulton--MacPherson boundary
compatibility. Equivalently, after descent it gives chain maps
\[
q_{g,n}\colon \cM_{\Ass}(g,n)_{\Sigma_n}
\longrightarrow \cM_{\Com}(g,n)
\]
compatible with the modular-operadic differentials and edge
compositions. With such a datum fixed, the \emph{averaging map}
\begin{equation}
\label{eq:averaging-map}
\av_R \;\colon\; \gEone \;\longrightarrow\; \gmod,
\qquad
\av_R(\phi)(g,n)
\;:=\;
\bigl[\phi(g,n)\bigr]_{R-\Sigma_n}
\end{equation}
is the induced $R$-twisted coinvariant comparison after the
descent identification. In the pole-free case $R=\id$, this is
the ordinary Reynolds quotient. Without this descent datum no map
\(\gEone\to\gmod\) is asserted.
\end{definition}

\begin{theorem}[Averaging is a dg Lie comparison on the descent locus]
\label{thm:av-dg-lie}
For an admissible averaging datum, the map
$\av_R \colon \gEone \to \gmod$ is a morphism of dg Lie algebras
on the descent locus. It commutes with the differential because
ribbon edge contraction descends to stable edge contraction under the
strong-unitary $R$-descent, and it commutes with the Lie bracket
because the twisted coinvariant comparison is compatible with
modular-operadic composition. Aritywise surjectivity holds after
choosing linear splittings of the finite descent quotients; no
canonical dg Lie section is part of the datum.
\end{theorem}

\begin{proof}
The hypotheses say exactly that the $R$-twisted descent quotient
intertwines edge-contraction differentials and modular-operadic
compositions. Applying $\Hom(-,\End_{\cA})$ and passing to the
descent quotient gives a chain map preserving the convolution
bracket. Over characteristic zero the finite group algebra
$k[\Sigma_n]$ is semisimple, so each arity admits linear splittings
of the quotient maps. These splittings are not canonical and need
not respect the full cross-arity differential.
\end{proof}

\begin{theorem}[Degree-$2$ averaging: Heisenberg exact, affine KM with Sugawara shift]
\label{thm:av-degree-2}
Let $r(z) \in \End(V \otimes V) \otimes \cO(*\Delta)$ be the
genus-$0$ degree-$2$ $E_{1}$ shadow of $\cA$: the degree-$2$
component of the $E_{1}$ Maurer--Cartan element, in the
pre-dualisation convention of
Volume~I (the shadow-archetype column of
\cite[Example]{Lorgat26I}, equivalently the collision residue
computed in $\cA\otimes\cA$ before applying the Koszul pairing).
Then
\begin{equation}
\label{eq:av-degree-2}
\av_{R,n=2}\bigl(r(z)\bigr)
\;=\;
\begin{cases}
\kappa(\cA), & \text{for abelian Heisenberg families},\\[4pt]
\kappa_{\mathrm{dp}}(\widehat\fg_{k})
:= \dfrac{k\,\dim(\fg)}{2h^{\vee}}, &
\text{for non-abelian affine Kac--Moody in the trace-form convention.}
\end{cases}
\end{equation}
In the affine Kac--Moody case the full modular characteristic is
\[
\kappa(\widehat\fg_{k})
\;=\;
\kappa_{\mathrm{dp}}(\widehat\fg_{k}) + \frac{\dim(\fg)}{2}
\;=\;
\frac{\dim(\fg)(k+h^{\vee})}{2h^{\vee}},
\]
so the extra $\dim(\fg)/2$ is the Sugawara simple-pole shift.
Concretely, for the
principal examples:
\begin{itemize}
\item Heisenberg $\cH_{k}$: $r(z) = k/z$, $\kappa(\cH_{k}) = k$.
\item Affine Kac--Moody $\widehat\fg_{k}$: $r(z) = k\Omega_{\fg}/z$
 (the level $k$ survives; at $k=0$ the $r$-matrix vanishes
 identically),
 $\av_R(r(z)) = \kappa_{\mathrm{dp}}(\widehat\fg_k)
 = k\dim(\fg)/(2h^{\vee})$, and
 $\kappa(\widehat\fg_{k}) = \dim(\fg)(k+h^{\vee})/(2h^{\vee})$.
\item Virasoro $\Vir_{c}$: $r(z) = (c/2)/z^{3} + 2T/z$ (cubic
 plus simple pole, no quartic), and $\kappa(\Vir_{c}) = c/2$.
\end{itemize}
\end{theorem}

\begin{proof}
The degree-$2$ Maurer--Cartan component is, by definition, the
degree-$2$ $E_{1}$ shadow
$r(z) \in \End(V^{\otimes 2}) \otimes \cO(*\Delta)$
in the pre-dualisation convention. Averaging over $\Sigma_{2}$
projects onto the trivial isotypic component. For abelian
Heisenberg families this scalar coincides with the full modular
characteristic. For non-abelian affine Kac--Moody in the
trace-form convention $r(z)=k\Omega_{\fg}/z$, the averaging sees
only the double-pole channel and yields
$\kappa_{\mathrm{dp}} = k\dim(\fg)/(2h^{\vee})$; the simple-pole
self-contraction contributes the additional Sugawara term
$\dim(\fg)/2$, so
$\kappa = \kappa_{\mathrm{dp}} + \dim(\fg)/2$. The example
computations are recorded in \cite[Chapter~5]{Lorgat26I}; each
formula is verified by the three independent paths of pole-order
analysis, the limiting case at $k=0$ for the trace-form
$r$-matrix, and cross-family consistency against the landscape
census.
\end{proof}

\begin{theorem}[Degree-$3$ averaging: associator to cubic shadow]
\label{thm:av-degree-3}
Let $\Phi_{\mathrm{KZ}}(\cA) \in \End(V^{\otimes 3})$ be the
genus-$0$ degree-$3$ Maurer--Cartan component, which for affine
Kac--Moody $\widehat\fg_{k}$ is the Drinfeld--Kontsevich KZ
associator. Then
\begin{equation}
\label{eq:av-degree-3}
\av_{R,n=3}\bigl(\Phi_{\mathrm{KZ}}(\cA)\bigr)
\;=\;
\mathfrak{C}(\cA),
\end{equation}
where $\mathfrak{C}(\cA)$ is the cubic shadow of Volume I,
equivalently the $\Sigma_{3}$-coinvariant of the associator. The
antisymmetric component $[\Omega_{12},\Omega_{23}]$ lies entirely
in $\ker(\av)$: it is antisymmetric under the transposition
$(1 \leftrightarrow 3)$, so its $\Sigma_{3}$-average vanishes.
Only the symmetric products
$\Omega_{12}\Omega_{23} + \Omega_{23}\Omega_{12}$
contribute to $\mathfrak{C}(\cA)$.
\end{theorem}

\begin{proof}
The $\Sigma_{3}$-isotypic decomposition of $\Phi_{\mathrm{KZ}}$ in
$\End(V^{\otimes 3})$ splits into the trivial, sign, and standard
representations. In the pole-free specialization, the Reynolds
operator $\av = (1/6)\sum_{\sigma
\in \Sigma_{3}}\sigma$ projects onto the trivial isotypic
component. A direct check on $[\Omega_{12},\Omega_{23}]$ shows
it is antisymmetric under $(1 \leftrightarrow 3)$ and therefore
has trivial isotypic component zero. The symmetric products
$\Omega_{ij}\Omega_{k\ell} + \Omega_{k\ell}\Omega_{ij}$ are the
only surviving contributors. The identification of this
coinvariant with the cubic shadow is Volume I Theorem D(iii).
\end{proof}

\begin{proposition}[Splittings and associator dependence]
\label{prop:nonsplitting}
For a fixed admissible averaging datum, the sequence of graded
linear spaces
\begin{equation}
\label{eq:nonsplitting-ses}
0 \longrightarrow \ker(\av) \longrightarrow \gEone
 \xrightarrow{\av} \gmod \longrightarrow 0
\end{equation}
splits aritywise, but it has no canonical splitting compatible with
the full Maurer--Cartan problem. At each fixed degree $n$, the
inclusion
$\End^{\Sigma_{n}}(V^{\otimes n}) \hookrightarrow
\End(V^{\otimes n})$ is a Lie subalgebra embedding. The obstruction
to a canonical global section is cross-degree: the bar differential
reduces degree by one and can carry a chosen representative in
\(\ker(\av_{R,n})\) to a term with nonzero
\(\av_{R,n-1}\)-component.
On the associator-controlled subproblem, choices of compatible
splitting form the usual Drinfeld associator torsor under
$\GRT_1$.
\end{proposition}

\begin{proof}
Maschke semisimplicity gives aritywise linear splittings over a
field of characteristic zero. Compatibility with the bar
differential is extra structure, because the differential relates
different arities. For KZ/Yangian input the remaining choice is the
choice of associator; Drinfeld's theorem identifies the set of
associators as a $\GRT_1$-torsor~\cite{Drinfeld90}, and
Etingof--Kazhdan quantisation supplies the corresponding
quasi-triangular deformation data~\cite{EtingofKazhdan96}.
\end{proof}

% ================================================================
\section{The five theorems as averaging-invariants}
\label{sec:five-theorems}

The following are conditional all-genus $E_1$ comparison surfaces.
Their genus-zero finite-degree pieces are proved where cited; the
all-genus statements assume the strict ordered Mittag--Leffler
tower, PBW-completeness on the Koszul locus, strong-unitary
$R$-descent, and completed finite-window convergence. In the
verified families, Theorem D identifies the scalar modular
characteristic with the degree-$2$ coinvariant of $r(z)$.

\begin{theorem}[Theorem~$\mathrm{A}^{E_{1}}$: ordered bar--cobar]
\label{thm:e1-theorem-a}
\textup{(}\emph{quadrant} Open, \emph{presentation} bar / twisting
coalgebra, \emph{level} $1 \leftrightarrow 2$ of the Beilinson
tower, \emph{hypothesis package}: strongly admissible augmented
cyclic $E_{1}$-chiral algebra $\cA$ with an admissible $R$-twisted
descent datum.\textup{)}
Let $\cA$ be a strongly admissible augmented cyclic
$E_1$-chiral algebra with an admissible descent datum. The genus-$g$
ordered bar complex
$\Barord(\cA)(g,n)$ is an $F\Ass$-coalgebra with differential
satisfying $d^{2} = 0$, where $F\Ass$ is the Feynman transform
of the ribbon modular operad. The bar and cobar functors form an
adjunction
\[
(\Cobar)^{(g)} \dashv (\Barord)^{(g)}
\colon \mathsf{Alg}_{F\Ass} \rightleftarrows
\mathsf{Coalg}_{F\Ass},
\]
and coinvariant projection recovers Theorem A of Volume I.
\end{theorem}

\begin{theorem}[Theorem~$\mathrm{B}^{E_{1}}$: ordered inversion on the Koszul locus]
\label{thm:e1-theorem-b}
\textup{(}\emph{quadrant} Open, \emph{presentation} bar / cobar
inversion, \emph{level} $1 \to 3$ of the Beilinson tower,
\emph{hypothesis package}: $E_{1}$ Koszul locus
\textup{(}PBW-completeness on the ordered bar
complex\textup{)}, propagated genus by genus via the MC
perturbation lemma.\textup{)}
Assume $\cA$ lies on the $E_{1}$ Koszul locus, meaning the
PBW-completeness hypothesis of
\cite[Chapter~4]{Lorgat26I} holds for the ordered bar complex.
Then at each genus $g$, the counit
\[
(\Cobar)^{(g)}\!\bigl(\Barord(\cA)(g,\bullet)\bigr)
\;\xrightarrow{\;\sim\;}\; \cA
\]
is a quasi-isomorphism (the inversion $\Cobar \circ \Barord
\simeq \id$ on the Koszul locus, not a Koszul-duality statement),
and the $E_{1}$ linear Koszul dual $\cA^{!,E_{1}}$ obtained by
linear duality on a finite-dimensional model of $\Barord(\cA)$ is
an $E_{1}$-chiral algebra at each such genus. The scope of ``all
genera'' here is contingent on the PBW-completeness hypothesis
being propagated genus by genus via the MC perturbation lemma of
\cite[Chapter~9]{Lorgat26I}.
\end{theorem}

\begin{theorem}[Theorem~$\mathrm{C}^{E_{1}}$: conditional ordered complementarity]
\label{thm:e1-theorem-c}
\textup{(}\emph{quadrant} Open, \emph{presentation} ordered
complementarity, \emph{level} $5$ per stratum, \emph{hypothesis
package}: ordered Koszul / descent locus, finite perfect model for
the ordered bar pairing, ordered Hochschild centre
$Z^{\mathrm{der},E_1}_{\mathrm{ch}}(\cA)$ identified with the
centre used by the strict ordered tower.\textup{)}
Assume $\cA$ lies in the ordered Koszul/descent locus, admits a
finite perfect model for the ordered bar pairing, and has an
ordered Hochschild centre
$Z^{\mathrm{der},E_1}_{\mathrm{ch}}(\cA)$ identified with the
centre used by the strict ordered tower. At each stable $(g,n)$
the ordered complementarity comparison is conditional:
\begin{equation}
\label{eq:e1-complementarity}
Q_{g,n}^{E_{1}}(\cA) + Q_{g,n}^{E_{1}}(\cA^{!,E_{1}})
\;\longrightarrow\;
H^{*}\!\bigl(\overline{\cM}_{g,n}^{\,\mathrm{rib}},
Z^{\mathrm{der},E_1}_{\mathrm{ch}}(\cA)\bigr),
\end{equation}
after applying the ordered centre functor and the perfect pairing.
Here the centre is a Hochschild object, not the subspace of
$R$-matrix-fixed vectors, and the comparison is not a consequence
of averaging alone. The scalar identity
$\kappa + \kappa^{!} = 0$ for Kac--Moody and free-field families
(and $\kappa + \kappa^{!} = 13$ for Virasoro) is the
$\Sigma_{n}$-coinvariant image in the verified scalar examples.
\end{theorem}

\begin{theorem}[Theorem~$\mathrm{D}^{E_{1}}$: degree-$2$ package]
\label{thm:e1-theorem-d}
\textup{(}\emph{quadrant} Open, \emph{presentation} degree-$2$
shadow tower, \emph{level} $4$ on $\overline\cM_{g,n}$,
\emph{hypothesis package}: genus-refined degree-$2$ projection of
$\Theta^{E_{1}}_\cA$ exists genus by genus; cyclic admissibility
on each stratum; higher-genus line-side $R$-matrix
interpretation requires additional Yangian input
\textup{(}open\textup{)}.\textup{)}
The formal ordered degree-$2$ shadow series
\begin{equation}
\label{eq:e1-theorem-d-series}
R^{E_{1},\mathrm{bin}}(z;\hbar)
\;=\;
\sum_{g \geq 0} \hbar^{2g}\, r_{g}(z)
\end{equation}
obtained from the genus-refined degree-$2$ projection of the
$E_{1}$ Maurer--Cartan element is the $E_{1}$ degree-$2$
characteristic package of $\cA$. It is universal, additive,
antisymmetric under opposite-duality, and satisfies
$\av_R(r_{0}(z)) = \kappa(\cA)$ on abelian families, while for
non-abelian affine Kac--Moody one has
$\av_R(r_{0}(z)) = \kappa_{\mathrm{dp}}(\cA)$ and
$\kappa(\cA) = \kappa_{\mathrm{dp}}(\cA) + \dim(\fg)/2$.
The coinvariant recovery of the scalar curvature from $r(z)$ is
the content of Theorem~\ref{thm:av-degree-2}. A full higher-genus
line-side modular $R$-matrix interpretation requires additional
Yangian input and remains open here.
\end{theorem}

\begin{theorem}[Theorem~$\mathrm{H}^{E_{1}}$: ordered chiral Hochschild]
\label{thm:e1-theorem-h}
\textup{(}\emph{quadrant} Open, \emph{presentation} ordered chiral
Hochschild concentration, \emph{level} $3$ of the Beilinson tower,
\emph{hypothesis package}: $\cA$ complete, ordered Koszul,
admissible $R$-twisted descent datum; complete $\cA$-bimodule $M$
in the same completion as $\cA$; chain-level statement requires
weight-completed / pro-object / $J$-adic ambient for class M.\textup{)}
Let $\cA$ be complete, ordered Koszul, and equipped with the
admissible descent datum of Definition~\ref{def:averaging-map}.
For complete bimodules $M$ in the same completion, the genus-$g$
ordered chiral Hochschild
homology is computed by the ordered chiral coHochschild complex of
the ordered bar coalgebra $C = \Barord(\cA)$. In the uniform-weight
case the twisted differential contains the curvature term
$\kappa(\cA)\lambda_g$. At genus $g\geq 2$ with multi-weight field
content the scalar formula receives the cross-channel correction
$\delta F_g^{\mathrm{cross}}$ of Volume I. The
$\Sigma_{n}$-coinvariant recovers the symmetric chiral Hochschild
complex of Theorem H.
\end{theorem}

\begin{remark}[Theorem D as the degree-two model]
\label{rem:theorem-d-degree-two-model}
Theorem D is the degree-two model for the $E_1$ primacy thesis.
The scalar $\kappa$ is one number per family, while its ordered
lift is the meromorphic collision function $r(z)$. For Heisenberg,
$\av_R(r(z))=\kappa$; for non-abelian affine Kac--Moody,
$\av_R(r(z))=\kappa_{\mathrm{dp}}$ and the Sugawara channel adds
$\dim(\fg)/2$. The projection from $r(z)$ to its scalar shadow
loses the braiding, quantum-group, and line-operator data.
\end{remark}

\begin{remark}[What averaging forgets]
\label{rem:what-av-forgets}
The kernel of $\av$ contains several objects of independent
interest. At degree $2$, $\ker(\av)$ is the antisymmetric part of
$r(z)$; for bosonic generators of even desuspended degree this
vanishes identically, and $\kappa$ recovers all of $r(z)$. For
generators of odd desuspended degree, the kernel is nontrivial
and captures the parity-dependent Eulerian weight structure
on the bar complex. At degree $3$, $\ker(\av)$ contains the
linearised Drinfeld commutator $[\Omega_{12},\Omega_{23}]$ and,
by the non-splitting proposition, obstructs a canonical
reconstruction of the full associator from $\mathfrak{C}$.
The $\GRT_{1}$-torsor attached to associator-controlled splittings
is the operadic shadow of Etingof--Kazhdan non-uniqueness of
quantisation.
\end{remark}

% ================================================================
\section{Three tiers of \texorpdfstring{$R$}{R}-matrix}
\label{sec:three-tiers}

The ordered bar complex distinguishes three tiers of input,
corresponding to how the degree-$2$ Maurer--Cartan component
$r(z)$ is constructed. The tiers are a refinement of the
$E_{n}$-hierarchy restricted to the degree-$2$ slice; only the
third is genuinely $E_{1}$ in the sense that the Maurer--Cartan
element cannot be recovered from the underlying
$E_{\infty}$-chiral algebra.

\subsection*{Tier (a): pole-free, $R$ is the Koszul-signed flip}

If $\cA$ is pole-free (that is, the OPE $a(z)b(w)$ has no
singularity for any $a, b \in \cA$), then the degree-$2$ collision
residue vanishes identically and $r(z) = 0$. In this tier the
ordered bar complex admits a canonical $\Sigma_{n}$-equivariant
structure: the descent datum $R_{\sigma}$ on the local system
$\mathcal{F}_{n}^{\mathrm{ord}}$ is the Koszul-signed permutation
\begin{equation}
\label{eq:tier-a-R}
R_{\sigma} \;=\; \varepsilon(\sigma)\cdot \tau_{\sigma},
\end{equation}
where $\tau_{\sigma}$ permutes tensor factors. The ordered and
symmetric bar complexes are quasi-isomorphic in this tier. Every
genuinely commutative (in the $E_{\infty}$ sense) example lives
here: the trivial chiral algebra, free-field tensor products
with vanishing cross-OPE, and the structure sheaf of a point.

\subsection*{Tier (b): vertex algebras, $R$ from local OPE}

If $\cA$ is a vertex algebra in the Borcherds--Frenkel--Lepowsky--Meurman sense (more generally, any $E_{\infty}$-chiral algebra with
poles in its OPE), then $r(z)$ is determined by the degree-$2$
residue of the OPE along the Arnold form $d\log(z_{1}-z_{2})$.
All standard families of conformal field theory live here:
Heisenberg $\cH_{k}$ with $r(z) = k/z$, affine Kac--Moody
$\widehat\fg_{k}$ with $r(z) = k\Omega_{\fg}/z$, Virasoro $\Vir_{c}$ with $r(z) = (c/2)/z^{3}
+ 2T/z$ (cubic and simple poles, no quartic), $\beta\gamma$,
$bc$, $W$-algebras, Bershadsky--Polyakov. In each case $r(z)$ is
fully determined by the underlying $E_{\infty}$ vertex algebra
structure; the $E_{1}$ refinement adds no new data at degree $2$
beyond what is already visible in the OPE. The averaging map on
tier (b) recovers the scalar degree-$2$ shadow: for Heisenberg
$\av_R(r(z)) = \kappa(\cA)$, while for non-abelian affine
Kac--Moody $\av_R(r(z)) = \kappa_{\mathrm{dp}}(\cA)$ and the full
$\kappa(\cA)$ adds the Sugawara shift $\dim(\fg)/2$; the kernel is
nontrivial but reconstructible from $\cA$.

\subsection*{Tier (c): genuinely $E_{1}$}

The third tier is populated by algebras in which $r(z)$ is
independent input: the Maurer--Cartan element cannot be recovered
from any underlying $E_{\infty}$ structure. The paradigmatic
examples are the Yangians $Y_{\hbar}(\fg)$ and the
Etingof--Kazhdan quantum vertex algebras~\cite{EtingofKazhdan00}
obtained from quasi-triangular Lie bialgebras. In this tier the
degree-$2$ Maurer--Cartan element is the rational Yang--Baxter
$r$-matrix $r_{\mathrm{YB}}(z)$, and the associator is a full
Drinfeld associator. Existing vertex-algebra theory cannot
recognise these algebras, and the ordered bar complex is the
minimal framework that makes them visible.

\begin{remark}[$E_{1}$ is the atom of the nonlocal regime]
\label{rem:e1-atom-nonlocal}
The three tiers form a strict hierarchy. Every standard chiral
algebra is $E_{\infty}$ in the sense that its OPE is symmetric in
the operator ordering up to the signs from the Arnold form; it is
the interpretation as an $E_{1}$ algebra that is the refinement.
Tier (c) consists of genuinely nonlocal objects, in the precise
sense that the Maurer--Cartan element lies outside the image of
any $E_{\infty}$-to-$E_{1}$ promotion functor. The error of
calling $E_{\infty}$-chiral algebras ``local'' and $E_{1}$-chiral
algebras ``non-abelian'' obscures the structure. The operative
distinction is symmetric versus ordered, or closed versus open.
The averaging map of Section~\ref{sec:av-map} erases the
ordering.
\end{remark}

% ================================================================
\section{Four functors, four distinct objects}
\label{sec:four-functors}

The structural part closes with a table of
functors. Four distinct functors act on a cyclic $E_{1}$-chiral
algebra $\cA$, and each produces a different object. Care is
required in keeping them apart.

\begin{enumerate}[label=\textup{(\roman*)}]
\item \emph{Bar.} The ordered bar functor
 $\cA \mapsto \Barord(\cA) = T^{c}(s^{-1}\bar\cA)$. The output
 is a coalgebra over the Feynman transform $F\Ass$ of the
 ribbon modular operad.
\item \emph{Cobar.} The cobar functor
 $C \mapsto \Cobar(C)$ on $F\Ass$-coalgebras. On the PBW-complete
 locus of Theorem~\ref{thm:e1-theorem-b}, the composite
 $\Cobar \circ \Barord$ is quasi-isomorphic to the identity; this
 is a statement about the bar--cobar adjunction returning $\cA$
 to itself, and is structurally distinct from the functor that
 produces the dual algebra.
\item \emph{Verdier.} The Verdier-duality functor
 $\cA \mapsto \cA^{!}_{\infty}$, where $\cA^{!}_{\infty}$ is the
 chain-level dual on the Ran space. Verdier duality intertwines
 the bar functor in the sense of~\cite[Theorem~A]{Lorgat26I};
 the strict linear Koszul dual $\cA^{!}$, obtained by taking a
 finite-dimensional model and applying linear duality, in
 general differs from $\cA^{!}_{\infty}$ and the two must be
 kept distinct.
\item \emph{Derived centre.} The chiral Hochschild cochain
 functor
 $\cA \mapsto Z^{\mathrm{der}}_{\mathrm{ch}}(\cA)
 = \mathrm{HH}^{*,\mathrm{ch}}(\cA)$.
 In the open-closed picture this plays the role of the bulk:
 the centre of the open (line-operator) algebra $\cA$ is the
 closed (bulk) sector. This functor is \emph{not} the bar--cobar
 composition; it is a separate construction, the subject of
 Theorem H.
\end{enumerate}

Theorem A of Volume I is the compatibility statement between (i)
and (iii): the bar functor intertwines with Verdier duality on
the Ran space. Theorem B (its $E_{1}$ version,
Theorem~\ref{thm:e1-theorem-b} above) is the internal statement
about (i) and (ii): on the PBW-complete locus the bar-cobar
composite returns $\cA$ up to quasi-isomorphism. Theorem H is the
statement about (iv): the derived centre coincides with the
ordered chiral Hochschild complex of the bar coalgebra. The four
objects are four different outputs; the statement
``$\Barord(\cA)$ computes the bulk'' is category-theoretically
ill-formed, because the closed sector is the derived centre, not the bar.

\medskip

In the three-pillar language of Volume I, the distinction is:
bar-cobar adjointness gives inversion; Verdier duality gives
$\cA^{!}$; chiral Hochschild cochains give the closed sector. The averaging map
$\av$ of Section~\ref{sec:av-map} intertwines the ordered and
symmetric versions of the bar functor on the descent locus. The
Verdier, cobar, and Hochschild functors require their own
compatibility data; they are not formal consequences of averaging
alone.

% ================================================================
\section{The Swiss-cheese / chiral Deligne--Tamarkin mechanism}
\label{sec:swiss-cheese-DT}

The bar functor and the averaging map both live on the open colour
of the chiral theory. The dimensional uplift to a
$3$-real-dimensional holomorphic--topological theory is not
accomplished by either of them: it is a property of the universal
Swiss-cheese pair built by the chiral Deligne--Tamarkin
construction. The bar coalgebra $\Barord(\cA)$ is the chain-level
computational model used to present that pair.

\begin{theorem}[Swiss-cheese / chiral Deligne--Tamarkin]
\label{thm:swiss-cheese-chiral-DT-N3}
\textup{(}\emph{quadrant} Open, \emph{presentation} two-coloured
chiral Swiss-cheese pair, \emph{level} $1 \to 3$ of the Beilinson
tower, \emph{hypothesis package}: a local $\Ainf$-chiral algebra
$\cA$ on a tangential log curve $(X, D, \tau)$ with completed
conformal-weight topology and finite chiral residue at every
double-point collision; a two-coloured chiral Swiss-cheese operad
$\SCchtop$ with chiral colour and topological collar
colour~\textup{\cite{Voronov99,KS00}}; an inner conformal vector
controlling $\bar Q$-translations on the topological collar.\textup{)}
The universal local chiral Swiss-cheese pair
\[
  U(\cA) \;=\; \bigl(\,
    \HH^{*,\mathrm{ch}}(\cA, \cA),\; \cA,\;
    \{\mu^{\mathrm{univ}}_{p;q}\}\,\bigr)
  \;=\;
  \bigl(\,Z^{\mathrm{der}}_{\mathrm{ch}}(\cA),\; \cA,\;
  \{\mu^{\mathrm{univ}}_{p;q}\}\,\bigr)
\]
is terminal in the category of local $\SCchtop$-pairs over $\cA$:
the closed colour $\HH^{*,\mathrm{ch}}(\cA, \cA)
= Z^{\mathrm{der}}_{\mathrm{ch}}(\cA)$ acts on the open colour
$\cA$ through the chiral Swiss-cheese mixed operations
$\mu_{p;q}$ \textup{(}$p$ closed inputs, $q$ open inputs\textup{)};
the open colour does not act back on the closed colour. For any
local $\SCchtop$-pair $(B, \cA, \{\mu_{p;q}\})$ over $\cA$ with
continuous mixed operations there is a unique morphism
$\Phi \colon (B, \cA, \mu) \to U(\cA)$ over the identity on $\cA$.
\end{theorem}

\begin{proof}[Citation]
This is Volume~I's chiral Deligne--Tamarkin theorem, the chiral
analogue of the Voronov Swiss-cheese formality
theorem~\cite{Voronov99} and the Kontsevich--Soibelman Deligne
conjecture~\cite{KS00}. Concretely, $\Phi(b)_n$ is the chiral
Hochschild cochain
$\Phi(b)_n(a_1, \ldots, a_n; \lambda_1, \ldots, \lambda_{n-1})
= (-1)^{|b|}\mu_{1;n}(b; a_1, \ldots, a_n; \lambda_1, \ldots,
\lambda_{n-1})$ on $n - 1$ open-sector spectral variables.
\end{proof}

\begin{remark}[The Swiss-cheese pair is the explanatory mechanism;
the bar is a computational model]
\label{rem:bar-is-computational-N3}
The $2$-real to $3$-real dimensional uplift is the universal
normal-collar direction of the codimension-one boundary condition
forced by the chiral Deligne--Tamarkin terminal-object construction
of Theorem~\textup{\ref{thm:swiss-cheese-chiral-DT-N3}}. It is not
produced by iterating the bar functor on $\cA$ alone. The bar
coalgebra $\Barord(\cA) = T^{c}(s^{-1}\bar\cA)$ is
$E_{1}$-coassociative on the curve geometry: a single
twisting coalgebra over the chiral associative operad
$(\mathrm{ChirAss})^!$, with deconcatenation coproduct and bar
differential constituting one graded coalgebra structure --- not
two colours of a Swiss-cheese algebra. The bar serves as a
chain-level computational model that resolves $\cA$ for the purpose
of computing the closed-colour Hochschild cochain complex via
\[
  Z^{\mathrm{der}}_{\mathrm{ch}}(\cA)
  \;\simeq\;
  \mathrm{Hom}_{\mathrm{ch}}\bigl(\Barord(\cA),\, \cA\bigr).
\]
The averaging map $\av_R$ of Section~\ref{sec:av-map} acts within
the open colour, projecting the ordered bar onto the symmetric bar;
it does not construct or witness the closed colour. The open and
closed colours are linked only through the Swiss-cheese mixed
operations $\mu_{p;q}$.

In particular, the slogans \emph{``$\Barord(\cA)$ is the bulk''},
\emph{``the bar direction explains the $2$d $\to$ $3$d HT lift''},
and \emph{``the bar is the open--closed transition''} are
category-level errors: the universal closed sector is the closed colour
$Z^{\mathrm{der}}_{\mathrm{ch}}(\cA)$, the dimensional uplift is
the universal collar of the Swiss-cheese pair, and the open--closed
transition is the chiral Deligne--Tamarkin morphism, not the bar
functor.
\end{remark}

\begin{remark}[Where the dimensional uplift comes from in
$\SCchtop$]
\label{rem:uplift-source-N3}
The chiral Swiss-cheese chain-level construction adapts the
topological Voronov pair~\cite{Voronov99} to the chiral setting:
the open colour records $E_1$-chiral operations on $\cA$ along the
curve $X$ (one complex direction, two real directions); the closed
colour records $E_2$-topological brace operations on
$Z^{\mathrm{der}}_{\mathrm{ch}}(\cA)$ in the transverse half-plane
(two real directions). The chiral Deligne--Tamarkin terminal-object
construction supplies a universal morphism between these two
sectors that is itself a chiral object, with the inner conformal
vector of the algebra controlling $\bar Q$-translations on the
topological collar. The total $\SCchtop$ structure is built from a
holomorphic chiral colour on the boundary curve and a topological
collar normal to it; the dimensional count is $2 + 1 = 3$. This
mechanism is neither the Costello--Gwilliam factorisation
construction, which is single-colour ambient, nor Voronov's
topological Swiss-cheese alone, which lacks a holomorphic colour
--- it is their fusion via chiral Deligne--Tamarkin. The bar of
$E_1$ is again $E_1$ (cofree conilpotent over $T^{c}$); the
Swiss-cheese pair $(E_2, E_1)$ at the chiral-topological level is
the source of the dimensional uplift.
\end{remark}

% ================================================================
\section{Conclusion and open problems}
\label{sec:open-problems}

The ordered bar complex $\Barord(\cA) =
T^{c}(s^{-1}\bar\cA)$, equipped with its deconcatenation
coproduct, is the level-$2$ chain object before averaging in modular
Koszul duality, and the symmetric bar $\BarSig(\cA)$ is its
$\Sigma_{n}$-coinvariant shadow under an admissible averaging
comparison $\av$. The five foundational theorems A, B, C, D, H of
modular Koszul duality admit ordered $E_{1}$ formulations on the
Koszul/descent locus. Theorem D is the degree-two calculation: the
scalar degree-$2$ shadow is recovered from the spectral
$r$-matrix $r(z)$ by averaging, giving $\kappa$
for abelian families and $\kappa_{\mathrm{dp}}$ for non-abelian
affine Kac--Moody before the Sugawara shift $\dim(\fg)/2$ is
added; the passage from $r(z)$ to the scalar shadow is a strict
projection with
nontrivial kernel.

Several problems remain open.

\subsection*{The full $E_{1}$ Koszul duality theorem}

Theorems~\ref{thm:e1-theorem-a} and~\ref{thm:e1-theorem-b} prove
the $F\Ass$-coalgebra structure, the inversion quasi-isomorphism on
the Koszul locus, and compatibility with averaging. They do not
prove an unconditional equivalence between all $E_1$-chiral
algebras and cocomplete $F\Ass$-coalgebras. A full equivalence of
categories requires tighter control of the coderived category of
$F\Ass$-coalgebras and a characterisation of the Koszul locus that
is independent of the commutative shadow.

\subsection*{The line-side modular $R$-matrix at higher genus}

Theorem~\ref{thm:e1-theorem-d} packages the ordered degree-$2$
shadow into a formal series
$R^{E_{1},\mathrm{bin}}(z;\hbar)$ whose $g$-th coefficient is the
ordered degree-$2$ correction at genus $g$. Whether this formal
series has a line-side interpretation as a genuine higher-genus
modular $R$-matrix is a separate question. Ribbon-graph physics
predicts such an object from the 't~Hooft expansion, and the cyclic
trace formalism gives the combinatorial language. Beyond genus one
this remains an open construction.

\subsection*{$\GRT_{1}$-torsor structure of splittings}

Proposition~\ref{prop:nonsplitting} isolates the point at which
aritywise linear splittings cease to be canonical: compatibility
with the Maurer--Cartan equation depends on associator data. The
Etingof--Kazhdan theorem provides this structure for Lie
bialgebras. The chiral analogue, including a functorial lift
$\kappa \mapsto r(z)$, remains open.

\subsection*{Swiss-cheese Koszul duality and the brane sector}

Theorem~\ref{thm:swiss-cheese-chiral-DT-N3} supplies the universal
local Swiss-cheese pair
$(Z^{\mathrm{der}}_{\mathrm{ch}}(\cA),\, \cA)$ as a terminal object,
identifying the bulk and the open--closed mechanism that lifts
$2$d chiral data to a $3$d holomorphic--topological theory. The
averaging map of this paper is internal to the open colour and
does not extend to a $\SCchtop$-functor. A Koszul duality theorem
for the chiral Swiss-cheese operad itself, that would package the
two-coloured operad together with its bar / cobar adjunction and
bind the brane / boundary-condition data of Volume~II to the
terminal pair, remains open. The mixed sector governing brane
boundaries and the open--closed coupling belongs to Volume~II.

\bigskip

In modular Koszul duality, the ordered bar complex $\Barord(\cA) =
T^{c}(s^{-1}\bar\cA)$ is the primitive object and the symmetric
bar is its $R$-twisted coinvariant shadow when descent exists. The
five theorem surfaces A, B, C, D, H are the invariants visible
after averaging. At degree $2$, Heisenberg satisfies
$\av_R(r(z)) = \kappa(\cA)$; non-abelian affine Kac--Moody satisfies
$\av_R(r(z)) = \kappa_{\mathrm{dp}}(\cA)$, with the Sugawara shift
$\dim(\fg)/2$ supplying the full $\kappa(\cA)$.

% ================================================================
\appendix
\section{Ordered-locus evidence}
\label{sec:ordered-locus-evidence}

The ordered comparison becomes a theorem on the locus where the
bar construction, the $R$-twisted descent, and the associator data
are all present.

\begin{definition}[Ordered Koszul chiral algebra]
\label{def:ordered-koszul-chiral}
A chiral algebra $\cA$ on a smooth curve $X$ is \emph{ordered
Koszul} if:
\begin{enumerate}[label=\textup{(\alph*)},leftmargin=2.2em]
\item the ordered bar complex
$\Barord(\cA)=T^c(s^{-1}\bar\cA)$ with deconcatenation coproduct
carries a coalgebra deformation retract onto its cohomology over
$(\mathrm{ChirAss})^!$;
\item the counit $\Cobar(\Barord(\cA))\to\cA$ is a
quasi-isomorphism in the Francis--Gaitsgory factorization ambient;
\item the $R$-twisted descent from ordered to symmetric bar is
compatible with the chiral twisting morphism.  This requires the
Yang--Baxter equation and the unitarity relation
$R(z)R^{\op}(-z)=\id$ after the scalar normalization appropriate
to the family.
\end{enumerate}
\end{definition}

\begin{proposition}[Yangian ordered Koszul instance]
\label{prop:yangian-ordered-koszul}
The chiral Yangian $Y_\hbar(\fg)^{\mathrm{ch}}$, defined from the
Drinfeld second presentation and Etingof--Kazhdan flatness, is
ordered Koszul in the sense of
Definition~\textup{\ref{def:ordered-koszul-chiral}}.
\end{proposition}

\begin{proof}
The Drinfeld second-presentation generators assemble an ordered bar
complex controlled by pure-braid Orlik--Solomon cohomology. The
Shelton--Yuzvinsky Koszulity theorem~\cite{SheltonYuzvinsky97}
gives the required quadratic
Koszul calculation on the pure-braid algebra, while
Etingof--Kazhdan flatness supplies the chiral twisting morphism and
the normalized $R$-twisted descent.
\end{proof}

\begin{theorem}[Ordered quantum-group comparison on the Koszul locus]
\label{thm:ordered-qg-comparison}
On the finite-type strongly admissible $E_1$ Koszul locus, after
fixing a Drinfeld associator $\Phi$ and a strong-unitary
$R$-twisted descent gauge, the following structures are equivalent
descriptions of the same ordered bar coalgebra:
\begin{enumerate}[label=\textup{(\roman*)},leftmargin=2.2em]
\item the classical $R$-matrix $r_\cA(z)$;
\item the $A_\infty$-structure transferred to $\Barord(\cA)$;
\item the deconcatenation coproduct on $\Barord(\cA)$ with its
compatible chiral twisting morphism, coassociative up to $\Phi$.
\end{enumerate}
The asserted $H^0$ associator-independence for the
$\mathfrak{sl}_2$ Yangian remains a separate first-order
computation until it is inscribed in the body theorem.
\end{theorem}

\begin{proof}
A choice of Drinfeld associator identifies the ordered tensor
calculus with the quasi-triangular deformation supplied by
Etingof--Kazhdan. The transferred $A_\infty$ operations are then
read from the same associator-controlled Maurer--Cartan element as
the $R$-matrix. Conversely, the twisting morphism and
deconcatenation coproduct recover the binary $R$-matrix and its
higher associator corrections.
\end{proof}

\begin{remark}[Further ordered-locus extensions]
Three extensions remain conditional.
For $\mathfrak{gl}_N$, the chiral Drinfeld double can be constructed
internally from Miura data, cross-arity Stasheff operations, and
$U(\widehat{\mathfrak{gl}}_1)$ screenings.  For
$\cW_\infty$, the ordered comparison belongs to the
weight-completed class-$\mathsf M$ category. In genus one, Felder's
dynamical $R(z,\tau)$ and the Fay trisecant identity
\cite{Felder94,Fay73} replace the pentagon identity; higher genus
requires theta-divisor input.
\end{remark}

\begin{remark}[Toroidal formal-disk comparison]
The toroidal formal-disk comparison uses Schiffmann--Vasserot
CoHA~\cite{SchiffmannVasserot13} and the Ding--Iohara--Miki
presentation~\cite{DingIohara97} of
$U_{q,t}(\ddot{\mathfrak{sl}}_N)$ with spectral parameters $(u,v)$.
It is a formal-disk statement. The global
$\mathbb{P}^1\times\mathbb{P}^1$ comparison requires the
class-$\mathsf{M}$ chain-level topologization of the original bar
complex.
\end{remark}

\begin{remark}[Degree-two constants]
The degree-two averaging constants are:
\begin{itemize}[leftmargin=1.8em]
\item Heisenberg $\cH_k$: $\av_R(r(z))=\kappa(\cH_k)=k$.
\item Non-abelian affine Kac--Moody $V_k(\fg)$:
$\av_R(r(z))=\kappa_{\mathrm{dp}}(V_k(\fg))
=k\dim(\fg)/(2h^\vee)$, and
$\kappa(V_k(\fg))=\dim(\fg)(k+h^\vee)/(2h^\vee)$ after the
Sugawara shift $\dim(\fg)/2$.
\item Virasoro $\Vir_c$:
$r(z)=(c/2)/z^3+2T/z$, and $\kappa(\Vir_c)=c/2$.
\item Principal $\cW_N$:
$\kappa(\cW_N)=c(H_N-1)$ with
$H_N=\sum_{j=1}^N 1/j$; for $N=2$ this recovers
$\kappa(\Vir_c)=c/2$.
\end{itemize}
\end{remark}

% ================================================================
\begin{thebibliography}{99}

\bibitem{BeilinsonDrinfeld04}
A.~Beilinson and V.~Drinfeld,
\emph{Chiral algebras},
American Mathematical Society Colloquium Publications, vol.~51,
AMS, Providence, RI, 2004.

\bibitem{Drinfeld90}
V.~G. Drinfeld,
\emph{On quasitriangular quasi-Hopf algebras and a group closely
connected with $\mathrm{Gal}(\overline{\mathbb{Q}}/\mathbb{Q})$},
Leningrad Math.~J. \textbf{2} (1991), 829--860.

\bibitem{EtingofKazhdan96}
P.~Etingof and D.~Kazhdan,
\emph{Quantization of Lie bialgebras, I},
Selecta Math. (N.S.) \textbf{2} (1996), 1--41.

\bibitem{EtingofKazhdan00}
P.~Etingof and D.~Kazhdan,
\emph{Quantization of Lie bialgebras, V: Quantum vertex operator
algebras},
Selecta Math. (N.S.) \textbf{6} (2000), 105--130.

\bibitem{FrancisGaitsgory12}
J.~Francis and D.~Gaitsgory,
\emph{Chiral Koszul duality},
Selecta Math. (N.S.) \textbf{18} (2012), 27--87.

\bibitem{Fay73}
J.~D. Fay,
\emph{Theta Functions on Riemann Surfaces},
Lecture Notes in Mathematics, vol.~352, Springer-Verlag,
Berlin-New York, 1973.

\bibitem{Felder94}
G.~Felder,
\emph{Conformal field theory and integrable systems associated to
elliptic curves},
in \emph{Proceedings of the International Congress of
Mathematicians (Z\"urich, 1994)}, Birkh\"auser, 1995,
1247--1255; arXiv:hep-th/9407154.

\bibitem{DingIohara97}
J.~Ding and K.~Iohara,
\emph{Generalization of Drinfeld quantum affine algebras},
Lett. Math. Phys. \textbf{41} (1997), no.~2, 181--193.

\bibitem{GetzlerJones94}
E.~Getzler and J.~D.~S. Jones,
\emph{Operads, homotopy algebra and iterated integrals for
double loop spaces},
preprint, 1994, arXiv:hep-th/9403055.

\bibitem{GetzlerKapranov98}
E.~Getzler and M.~M. Kapranov,
\emph{Modular operads},
Compositio Math.\ \textbf{110} (1998), 65--126.

\bibitem{LodayVallette12}
J.-L. Loday and B.~Vallette,
\emph{Algebraic operads},
Grundlehren der mathematischen Wissenschaften, vol.~346,
Springer, Heidelberg, 2012.

\bibitem{Lorgat26I}
R.~Lorgat,
\emph{Modular Koszul Duality, Volume I},
monograph, 2026.

\bibitem{SchiffmannVasserot13}
O.~Schiffmann and E.~Vasserot,
\emph{Cherednik algebras, W-algebras and the equivariant
cohomology of the moduli space of instantons on
$\mathbb{A}^2$},
Publ. Math. IH\'ES \textbf{118} (2013), 213--342,
arXiv:1202.2756.

\bibitem{SheltonYuzvinsky97}
B.~Shelton and S.~Yuzvinsky,
\emph{Koszul algebras from graphs and hyperplane arrangements},
J.\ London Math.\ Soc.\ \textbf{56} (1997), 477--490.

\bibitem{Stasheff63}
J.~D. Stasheff,
\emph{Homotopy associativity of $H$-spaces, I},
Trans.\ Amer.\ Math.\ Soc.\ \textbf{108} (1963), 275--292.

\bibitem{Vallette14}
B.~Vallette,
\emph{Homotopy theory of homotopy algebras},
Ann.\ Inst.\ Fourier \textbf{70} (2020), 683--738.

\bibitem{Voronov99}
A.\,A.~Voronov,
\emph{The Swiss-cheese operad},
Contemp.\ Math.\ \textbf{239} (1999), 365--373;
arXiv:math/9807037.

\bibitem{KS00}
M.~Kontsevich and Y.~Soibelman,
\emph{Deformations of algebras over operads and the Deligne
conjecture},
in \emph{Conf\'erence Mosh\'e Flato 1999, Vol.~I}, Math.\ Phys.\
Stud.\ \textbf{21}, Kluwer, 2000, 255--307;
arXiv:math/0001151.

\end{thebibliography}

\end{document}
